\chapter{Prácticas guiadas}
\label{cap:practicas}

Las prácticas están ordenadas de menor a mayor dificultad. Cada una indica sus objetivos, el procedimiento con el código necesario, las cuestiones a responder y lo que debe entregarse. Al empezar cada sesión conviene ejecutar:
\begin{lstlisting}
clear; close all; cp = CartPole(); cp.clearSignals();
\end{lstlisting}
Los valores numéricos que se citan como referencia corresponden a los parámetros por defecto.

% ----------------------------------------------------------------
\stepcounter{practica}
\begin{practica}{Exploración en lazo abierto y validación del modelo}
\objetivos
\begin{itemize}
  \item Familiarizarse con la interfaz y comprobar el modelo frente a la teoría.
  \item Medir la frecuencia y el amortiguamiento de la oscilación colgante.
\end{itemize}
\procedimiento
\begin{enumerate}
  \item Ensayo de caída libre desde arriba (sin control):
\begin{lstlisting}
cp.setMode('open'); cp.setInitialState([0 0 0.01 0]);
res = cp.run(3); cp_plot_results(res); cp.animate(res);
\end{lstlisting}
  \item Ajustar a los primeros \SI{0.5}{s} una exponencial $\theta(t)\approx\theta_0\cosh(pt)$ y estimar $p$.
  \item Oscilación libre colgando: \code{cp.setInitialState([0 0 pi-0.2 0]); res = cp.run(20);}. Medir el periodo y el decremento logarítmico $\delta=\ln(A_k/A_{k+1})$, y calcular $\zeta=\delta/\sqrt{4\pi^2+\delta^2}$.
  \item Repetir el punto anterior con el carro bloqueado (\code{cp.setParam('M',1e3)}) y comparar con $\omega_0=\sqrt{mgl/\Jp}$.
\end{enumerate}
\preguntas
\begin{enumerate}
  \item ¿Coincide $p$ con el polo inestable de la tabla~\ref{tab:polos} ($5{,}45$)? ¿Por qué no exactamente?
  \item Explicar por qué la frecuencia con el carro libre (\SI{0.89}{Hz}) es mayor que con el carro bloqueado (\SI{0.79}{Hz}).
  \item ¿Qué fricción domina el amortiguamiento de la oscilación, la del carro o la del pivote? Comprobarlo variando \code{bc} y \code{bp}.
\end{enumerate}
\entregables Gráficas de los tres ensayos y una tabla comparativa entre teoría y simulación.
\end{practica}

% ----------------------------------------------------------------
\stepcounter{practica}
\begin{practica}{Identificación experimental con una señal chirp}
\objetivos Obtener la respuesta en frecuencia del sistema colgante a partir de datos y compararla con la figura~\ref{fig:bode}.
\procedimiento
\begin{enumerate}
  \item Generar un chirp de fuerza de \SI{0.7}{N} entre \SI{0.5}{Hz} y \SI{3}{Hz} durante \SI{20}{s} y aplicarlo en lazo abierto:
\begin{lstlisting}
T = 20; tt = (0:cp.P.h:T)'; f0 = 0.5; f1 = 3;
u = 0.7*sin(2*pi*(f0*tt + (f1-f0)/(2*T)*tt.^2));
cp.setInitialState([0 0 pi 0]); cp.setMode('open');
cp.setForceSignal(tt, u); res = cp.run(T); cp.clearSignals();
\end{lstlisting}
  \item Estimar la respuesta en frecuencia con \code{tfestimate(res.Fact, mod(res.th,2*pi)-pi, [], [], [], 1/cp.P.h)}, de la Signal Processing Toolbox, o dividiendo directamente las FFT.
  \item Ajustar un modelo de segundo orden a $\Theta/F$ cerca de la resonancia.
\end{enumerate}
\preguntas
\begin{enumerate}
  \item ¿Dónde aparece la resonancia y dónde la antirresonancia de $X/F$?
  \item ¿Por qué no conviene excitar por debajo de \SI{0.3}{Hz} en lazo abierto? (Pista: polo en el origen y longitud del raíl.)
  \item ¿Cómo afecta la fricción de Coulomb a la identificación cuando la amplitud de la fuerza es pequeña?
\end{enumerate}
\entregables Bode experimental superpuesto al teórico y parámetros identificados.
\end{practica}

% ----------------------------------------------------------------
\stepcounter{practica}
\begin{practica}{Diseño de un LQR: influencia de $Q$ y $R$}
\objetivos Entender el compromiso entre rapidez, esfuerzo y recorrido del carro.
\procedimiento
\begin{enumerate}
  \item Desde $\theta_0=\SI{0.15}{rad}$, repetir el ensayo con $Q_\theta\in\{20,200,2000\}$ y luego con $R\in\{0{,}01,\ 0{,}1,\ 1\}$ (véase la figura~\ref{fig:qsweep}).
  \item Para cada diseño, anotar $K$, los polos del lazo cerrado (\code{cp.Lin.up.poles\_CL}), el tiempo de asentamiento a $0{,}5^\circ$, $F_{\max}$ y $|x|_{\max}$.
  \item Aplicar la regla de Bryson con $\theta_{\max}=5^\circ$, $x_{\max}=\SI{0.2}{m}$ y $F_{\max}=\SI{10}{N}$. Comparar el resultado con el diseño por defecto.
  \item En la GUI, cambiar $Q$ y $R$ en vivo mientras se aplican empujones.
\end{enumerate}
\preguntas
\begin{enumerate}
  \item Comprobar que $K_x=-\sqrt{Q_x/R}$ en todos los diseños y justificarlo.
  \item ¿Qué ocurre cuando el diseño satura el actuador? ¿Hay un ángulo inicial máximo recuperable? Encontrarlo por bisección.
\end{enumerate}
\entregables Tabla de diseños con sus métricas y una recomendación justificada.
\end{practica}

% ----------------------------------------------------------------
\stepcounter{practica}
\begin{practica}{Seguimiento de referencias y fase no mínima}
\objetivos Observar la respuesta inversa y cuantificar los límites de seguimiento.
\procedimiento
\begin{enumerate}
  \item Aplicar un escalón con \code{cp.setRefSignal([0 1 1.001 8], [0 0 0.3 0.3])} y medir el retroceso inicial (figura~\ref{fig:lqrstep}).
  \item Referencias senoidales de \SI{0.2}{m} a 0,05; 0,1; 0,2 y \SI{0.4}{Hz}: medir la ganancia y el desfase entre $x$ y $r$, y dibujar el Bode del lazo cerrado.
  \item Implementar el modo 5 con acción integral (sección~\ref{sec:nuevo_modo}) y repetir el ensayo de escalón con la fricción de Coulomb aumentada a \code{Fc = 1}.
\end{enumerate}
\preguntas
\begin{enumerate}
  \item Relacionar el ancho de banda medido con la posición del cero $z=4{,}91$.
  \item ¿Desaparece la respuesta inversa con el integrador? ¿Por qué no puede desaparecer con ningún controlador?
  \item Proponer un prefiltro o feedforward que reduzca el retraso con referencias senoidales.
\end{enumerate}
\entregables Gráficas, Bode del lazo cerrado y comparación con y sin integral.
\end{practica}

% ----------------------------------------------------------------
\stepcounter{practica}
\begin{practica}{Rechazo de perturbaciones y robustez}
\objetivos Evaluar el LQR frente a empujones y a errores en los parámetros.
\procedimiento
\begin{enumerate}
  \item Programar empujones de amplitud creciente (\SI{0.5}{N} a \SI{3}{N}, durante \SI{0.1}{s}) con \code{setDisturbanceSignal} y determinar la amplitud máxima admisible antes de que salte la seguridad.
  \item Sin recalcular $K$, variar la masa real del carro, la masa de la varilla y su longitud. Para ello, usar \code{cp.setParamRaw('M', 1)}, que cambia la planta sin recalcular $K$, en lugar de \code{setParam}, que sí lo recalcularía.
  \item Repetir con el ruido de los sensores activado (\code{cp.setParam('noise\_on',1)}).
\end{enumerate}
\preguntas
\begin{enumerate}
  \item ¿Cuál de los parámetros es más crítico? Relacionarlo con la expresión de $A$ en~\eqref{eq:AB}.
  \item ¿Qué margen de ganancia tiene el lazo real? Estimarlo multiplicando $K$ por un factor $\alpha$ hasta que el sistema se vuelva inestable, y comparar con el margen teórico del LQR.
\end{enumerate}
\entregables Mapa de estabilidad en función del parámetro y de $\alpha$.
\end{practica}

% ----------------------------------------------------------------
\stepcounter{practica}
\begin{practica}{PID+PD frente a LQR}
\objetivos Ajustar el controlador clásico y compararlo con el óptimo.
\procedimiento
\begin{enumerate}
  \item Con el modo \code{'pid'}, empezar con $K_{px}=K_{dx}=0$ y ajustar $K_{p\theta}$ y $K_{d\theta}$ hasta estabilizar el ángulo. Observar cómo deriva el carro.
  \item Añadir $K_{px}$ y $K_{dx}$, primero con signo negativo y luego con signo positivo. Explicar lo que ocurre.
  \item Calcular los polos de $A-BK_{\text{PID}}$ para cada ajuste y comparar con los del LQR.
  \item Activar $K_{i\theta}$ y aplicar una fuerza constante de \SI{1}{N} (\code{cp.setForce(1)}). ¿Qué hace el integrador?
\end{enumerate}
\preguntas
\begin{enumerate}
  \item ¿Puede un ajuste PID reproducir exactamente el LQR? Razonar sobre la estructura de~\eqref{eq:pid}.
  \item Ventajas prácticas del PID cuando el modelo es poco conocido.
\end{enumerate}
\entregables Ajuste final, polos y ensayo comparativo.
\end{practica}

% ----------------------------------------------------------------
\stepcounter{practica}
\begin{practica}{\emph{Swing-up} por energía}
\objetivos Entender el control no lineal basado en energía y la conmutación entre controladores.
\procedimiento
\begin{enumerate}
  \item Ejecutar el \emph{swing-up} por defecto (figura~\ref{fig:swingup}) y verificar que $E(t)$ crece de forma monótona hasta la captura.
  \item Barrer $k_E\in\{2,5,10,20\}$ y $a_{\max}\in\{2,3,4,6\}$. Registrar el tiempo de captura y $|x|_{\max}$, y marcar los casos que activan la seguridad.
  \item Eliminar el recentrado (\code{kx\_sw = 0}, \code{kv\_sw = 0}) y observar la deriva del carro.
  \item Reducir \code{th\_catch} a \SI{0.15}{rad} y aumentarlo a \SI{0.6}{rad}. ¿Qué limita la región de captura?
\end{enumerate}
\preguntas
\begin{enumerate}
  \item Demostrar que $\dot V\le0$ con la ley~\eqref{eq:swing} sin saturación ni recentrado.
  \item ¿Por qué no basta con un LQR con mucha ganancia para levantar la varilla?
  \item Proponer una forma de reducir el recorrido del carro, por ejemplo con una energía objetivo ligeramente mayor que cero o con bombeo en varias etapas.
\end{enumerate}
\entregables Mapa (tiempo de captura y recorrido) frente a $(k_E,a_{\max})$ y explicación de la conmutación.
\end{practica}

% ----------------------------------------------------------------
\stepcounter{practica}
\begin{practica}{Grúa con anti-balanceo}
\objetivos Diseñar un posicionamiento rápido sin oscilación residual de la carga.
\procedimiento
\begin{enumerate}
  \item Reproducir la figura~\ref{fig:grua}: el modo \code{'crane'} con referencia cuadrada (\code{cp.setRef('square',0,0.4,0.1)}), frente al mismo modo con las ganancias angulares a cero:
\begin{lstlisting}
Kn = cp.P.K_dn; Kn(3:4) = 0;
cp.setParamRaw('K_dn', Kn);        % sin recalcular el LQR
\end{lstlisting}
  \item Diseñar un \emph{input shaper} ZV (dos impulsos separados medio periodo de la oscilación) para la referencia del modo PD y compararlo con el LQR.
  \item Variar la longitud de la varilla $\pm30\,\%$ sin rediseñar. ¿Qué método es más robusto?
\end{enumerate}
\preguntas Comparar tiempo de posicionamiento, balanceo residual y sensibilidad a la longitud de la varilla.
\end{practica}

% ----------------------------------------------------------------
\stepcounter{practica}
\begin{practica}{Efecto de la instrumentación}
\objetivos Cuantificar cómo afectan el muestreo, la cuantización, el ruido y el filtro derivativo.
\procedimiento
\begin{enumerate}
  \item Barrer $T_s\in\{1,2,5,10,20\}$\,ms. Mantener $h\le T_s/2$ con \code{cp.setParam('h',min(1e-3,Ts/2))}; \code{run()} ajusta \code{FixedStep} automáticamente. ¿A partir de qué $T_s$ se degrada el lazo? ¿Y cuándo se vuelve inestable?
  \item Reducir la resolución del encoder angular a 1024, 512 y 256 cuentas por vuelta y observar el «ruido» en $F$.
  \item Barrer $\tau_v\in\{1,5,10,30,100\}$\,ms. Representar la varianza de $F$ en régimen permanente frente al tiempo de asentamiento.
\end{enumerate}
\preguntas
\begin{enumerate}
  \item Explicar el compromiso de $\tau_v$ (figura~\ref{fig:velest}).
  \item ¿Qué resolución mínima de encoder recomendaríais para el banco real?
\end{enumerate}
\end{practica}

% ----------------------------------------------------------------
\stepcounter{practica}
\begin{practica}{Proyecto abierto: observador de estado}
\objetivos Sustituir la derivada filtrada por un observador basado en el modelo.
\procedimiento
\begin{enumerate}
  \item Discretizar el modelo lineal: $A_d=e^{AT_s}$, $B_d=\int_0^{T_s}e^{A\tau}B\,\dd\tau$ (véase la sección~7 de \code{examples.m}).
  \item Diseñar un observador de Luenberger $\hat{\vect s}_{k+1}=A_d\hat{\vect s}_k+B_dF_k+L(\vect y_k-C\hat{\vect s}_k)$ con polos 3--5 veces más rápidos que el lazo, o un filtro de Kalman estacionario tomando como varianzas de medida $q^2/12+\sigma^2$.
  \item Implementarlo como modo 6 en \code{cp\_controller\_block.m}, con \code{persistent shat}, y comparar con el modo 1 con ruido activado.
  \item Probar a realimentar solo $x$ (observable) o solo $\theta$ (no observable).
\end{enumerate}
\entregables Código del modo nuevo, comparación de la varianza de $F$ y del tiempo de asentamiento, y discusión del principio de separación.
\end{practica}
